%% Hand-authored circuitikz figure -- NOT generated by maestro2tex.
%% The cell is drawn as the abstracted symbol of sch_cell_symbol.tex, the same
%% circle-with-three-leads used in the Virtuoso schematics. The internal network
%% is deliberately NOT shown: the element values in note_models.tex define the
%% model. Do not restate them here.
\providecommand{\MaestroRoot}{include/analog/}
\input{\MaestroRoot sch_cell_symbol.tex}
\begin{figure}[htbp]
  \centering
\begin{circuitikz}[
    every node/.append style={font=\footnotesize},
    nlab/.style={font=\scriptsize, inner sep=1.5pt},
  ]
\CellSymThree{0,0}{1.0}
\end{circuitikz}
  \caption[{Three-electrode cell model.}]{Three-electrode cell model, the load the potentiostat actually drives, shown as the symbol used throughout this chapter. It is a lumped linear network between the counter, reference and working electrodes; its elements and their values are given in Table~\ref{tab:analog-models}. The model is frequency-independent (no Warburg diffusion term and no potential dependence), so it represents small-signal behaviour about a bias point rather than a full electrochemical response.}
  \label{fig:randles-cell}
\end{figure}
